\section{Заключение}
\label{sec:conclusion}

В работе исследовано влияние выбора реализации матричного умножения в
функциональном симуляторе на вычислительную точность при инференсе и обучении
нейросетей. Получены следующие основные результаты.

\begin{enumerate}\setlength\itemsep{2pt}
    \item \textbf{Модель распространения расхождений.} Построена аналитическая
        модель накопления аппаратных расхождений через слои сети при прямом и
        обратном проходах на основе линейного приближения и матриц Якоби
        (раздел~\ref{sec:model}); установлена роль слоёв с $\|J\|_2<1$
        (нормализация, насыщающиеся нелинейности, усреднения) как стоков
        расхождений (утверждение~\ref{st:sink}).
    \item \textbf{Методология и инструментарий.} Разработан измерительный
        конвейер, раздельно оценивающий вклад компилятора и реализации GEMM и
        классифицирующий природу ошибки за один прямой проход
        (раздел~\ref{sec:method}); он применим к сравнению любых пар реализаций
        вычислительных блоков, а не только GEMM.
    \item \textbf{Приемлемость GPU-реализации.} Доминирующим источником
        расхождений является компилятор XLA
        ($d_{\text{rel}}\approx2.2\cdot10^{-3}$), тогда как вклад GPU-реализаций
        GEMM на 2--5 порядков меньше (рис.~\ref{fig:hierarchy}). Для случайной
        ошибки коэффициент накопления остаётся умеренным ($R_{\text{fwd}}\le1.7$)
        во всех исследованных архитектурах. Следовательно, \emph{относительно
        принятой CPU-реализации как точки отсчёта} перенос GEMM на GPU не ухудшает
        вычислительную точность симулятора.
    \item \textbf{Природа ошибки важнее её нормы для накопления.} Поэлементная
        норма расхождения не является достаточным критерием выбора реализации:
        метрика смещения $b$ оказывается более информативным предиктором
        накопления, чем норма $d_{\text{rel}}$. Так, смещённая ошибка cuBLAS
        ($b\approx0.28$) при норме на 2--3 порядка меньшей, чем у случайной
        ошибки CUDA-ядра, формально даёт более высокий коэффициент накопления.
        Данный вывод согласуется с теорией конструктивного суммирования смещённых
        ошибок~\eqref{eq:scaling}, однако послойный сигнал cuBLAS лежит на уровне
        машинного эпсилона \texttt{FP32}, поэтому количественные значения
        $R$ для этой пары следует рассматривать как указание, а не как строго
        установленный эффект; его прямая проверка требует измерений повышенной
        разрядности.
\end{enumerate}

Среди двух GPU-реализаций cuBLAS меньше изменяет поведение симулятора
относительно CPU-эталона, но вносит систематическое смещение (с оговоркой о
надёжности измерений на уровне $\varepsilon_{\texttt{FP32}}$); CUDA-ядро не
создаёт смещения и даёт расхождение, надёжно измеримое в \texttt{FP32}. Все
выводы сформулированы относительно CPU-эталона; близость к арифметике конкретного
устройства требует отдельной верификации на реальном оборудовании.

Практическая значимость состоит в том, что предложенная методология позволяет
обоснованно выбирать реализацию GEMM в функциональном симуляторе без
исчерпывающего регрессионного тестирования: достаточно одного прямого прохода с
диагностикой типа ошибки. Перспективными направлениями представляются локализация
компиляторного вклада по механизмам (слияние ядер, выбор алгоритма свёртки,
неявная смена формата), оценка влияния накопленных расхождений на сходимость
многошагового обучения, распространение анализа на рекуррентные и
авторегрессивные архитектуры и верификация арифметической близости реализаций к
реальному устройству.
